\chapter*{Companion Project Files}
\addcontentsline{toc}{chapter}{Companion Project Files}

The source distribution is organised so that the manuscript, executable
notebooks and reusable Python package can evolve together.

\begin{verbatim}
introduction-to-image-processing-with-python/
|-- book/
|   |-- main.tex
|   |-- preamble.tex
|   |-- references.bib
|   |-- frontmatter/
|   |-- chapters/
|   |-- appendices/
|   `-- figures/
|-- notebooks/
|-- project/
|   |-- pyproject.toml
|   |-- src/pyimage/
|   |-- tests/
|   |-- examples/
|   `-- images/
`-- source-notes/
\end{verbatim}

The notebooks are the executable companion to individual chapters. The
\pyimage{} package is the finished reusable implementation. The test suite is
kept outside the notebooks so that the complete package can be verified from
the command line or from the integrated terminal in Visual Studio Code.

Publication releases should eventually tag the companion project with a
version corresponding to the edition of the book.
